\documentclass[12pt,a4paper]{article}
\usepackage[utf8]{inputenc}
\usepackage[T2A]{fontenc}
\usepackage[russian]{babel}
\usepackage{amsmath,amssymb,amsthm}
\usepackage{booktabs,array}
\usepackage{microtype}
\usepackage{geometry}
\geometry{margin=2.3cm}
\setlength{\emergencystretch}{2em}

\newtheorem{proposition}{Предложение}
\newtheorem{nogocriterion}{Критерий провала}

\title{Единый parent trace для gauge- и family-секторов}
\author{}
\date{4 августа 2026 г.}

\begin{document}
\maketitle

\section{Общий carrier}

Объединим affine spin-menu и одно anomaly-free \(SU(5)\)-поколение:
\[
\mathcal H_{parent}
=\mathbb C^4_{menu}
\otimes
\left(\mathbf{10}\oplus\overline{\mathbf5}\right).
\]
Его размерность равна
\[
4\cdot15=60,
\]
а полная operator algebra есть
\[
\mathcal A_{parent}=M_{60}(\mathbb C).
\]
Полная матричная algebra имеет единственный normalized trace
\[
\tau_{60}(X)=\frac1{60}\operatorname{Tr}_{60}X.
\]
Следовательно, central relative weights отсутствуют.

\section{Физический family sector}

Canonical affine-triplet projector равен
\[
P_{phys}=P_3\otimes I_{15},
\qquad
\operatorname{rank}P_{phys}=3\cdot15=45.
\]
Projected algebra изоморфна \(M_{45}(\mathbb C)\) и снова имеет единственный normalized conditional trace
\[
\tau_{phys}(X)
=\frac1{45}\operatorname{Tr}(P_{phys}XP_{phys}).
\]

\section{Family operator в том же trace}

Геометрический odd projector задаёт
\[
P_{heavy}=P_-\otimes I_{15}.
\]
Поскольку \(P_-\leq P_3\), имеем
\[
P_{phys}P_{heavy}=P_{heavy}.
\]
Его ранг равен
\[
\operatorname{rank}P_{heavy}=1\cdot15=15.
\]
Поэтому один parent trace фиксирует conditional family weight
\[
\tau_{phys}(P_{heavy})
=\frac{15}{45}
=\frac13.
\]

Это число не является коэффициентом старой \(\tau\)-формулы. Оно имеет другой и строгий смысл: доля heavy-family channel внутри трёхпоколенческого matter sector.

\section{Gauge trace}

Для одного \(\mathbf{10}\oplus\overline{\mathbf5}\) package нормированные subgroup indices равны
\[
I_{SU(3)}=I_{SU(2)}=I_{U(1)}=2.
\]
Family-triplet projector умножает каждый loop trace на один и тот же rank:
\[
I_{SU(3)}^{(3fam)}
=I_{SU(2)}^{(3fam)}
=I_{U(1)}^{(3fam)}
=3\cdot2=6.
\]

\begin{proposition}[Совместимость gauge и family reductions]
Rank-one family operator и равные \(SU(5)\)-нормированные gauge indices получаются из одного matrix trace на \(M_{60}(\mathbb C)\) и его единственного conditional trace на physical rank-45 sector. Независимый relative trace coefficient не требуется.
\end{proposition}

\section{Что именно закрыто}

Вторая сумасбродная гипотеза теперь имеет конкретное положительное содержание. Разные эффективные меры действительно могут быть тенями одного parent state:
\begin{itemize}
    \item gauge loop суммирует все три family copies;
    \item family mass operator использует spectrum rank-one projector;
    \item нормированная particle state имеет единичную норму;
    \item все эти операции совместимы с одним trace, а не выбираются независимо.
\end{itemize}

\section{Что не закрыто}

Trace фиксирует относительные multiplicities, но не задаёт action coefficients. В частности, он не определяет:
\begin{itemize}
    \item масштаб heavy-family eigenvalue;
    \item Higgs representation и vacuum expectation;
    \item коэффициент geometric shear term относительно gauge kinetic term;
    \item массы двух light families и mixing.
\end{itemize}

\begin{nogocriterion}[Action-level gate]
Алгебраическое единство trace не считается единством физического действия, пока наиболее общий \(SU(5)\times AGL(2,2)\)-инвариантный action не фиксирует относительный коэффициент family-breaking term. Если этот коэффициент свободен, гипотеза остаётся правильным bookkeeping, но не предсказательной динамикой.
\end{nogocriterion}

\section{Вердикт}

Вторая идея оказалась содержательнее простой фразы о разных мерах. В minimal tensor-product model parent trace действительно единственен и одновременно поддерживает:
\[
\text{три anomaly-free поколения},
\qquad
\text{равные gauge indices},
\qquad
\text{rank-one family channel}.
\]
Следующий решающий вопрос снова относится к действию, а не к числам: фиксирует ли symmetry относительную жёсткость gauge- и family-направлений или оставляет новую непрерывную ручку.

\end{document}